% Root and D full literal audits passed, including all derived shifts and module sides.
\subsection{Two-column graph actions and their contracted algebra}

\begin{lemma}[A two-term complex over a local algebra]
\label{C:local-two-term-derived-tensor}
Let $A$ be a finite-dimensional local rational algebra with
radical $J$ and residue field $A/J=\mathbb Q$.  Let $D$ be a
complex of finite-dimensional right $A$-modules concentrated in
homological degrees zero and one.  If
$H_i(D\otimes_A^{\mathbb L}\mathbb Q)$ vanishes for all
sufficiently large $i$, then it vanishes for every $i\geq2$.
\end{lemma}
\begin{proof}
Choose a bounded-below minimal resolution of $D$ by finitely
generated free right $A$-modules, in nonnegative degrees.
One construction starts with a bounded-below free resolution
and successively splits off each invertible matrix entry as a
contractible two-term free summand.  Over a local algebra the
remaining differentials have all entries in $J$; each degree
requires only finitely many cancellations.  This is the usual
minimal projective-complex construction, applied to the given
two-term complex.

Tensoring that minimal complex with the left residue module
$\mathbb Q$ makes all differentials zero.  The stated eventual
vanishing and Nakayama's lemma therefore make the free resolution
bounded.  Let $N$ be its largest nonzero degree.  If $N>1$, the
vanishing of $H_N(D)$ forces the top differential to be injective.
This is impossible.  A matrix of a right-module map has entries
acting on the left of the input coefficients.  Since $J$ is
nilpotent, there is a nonzero $a\in A$ with $Ja=0$; take a
nonzero element of the last nonzero radical power, or $a=1$ if
$J=0$.  The vector which is $a$ in one top free coordinate and
zero elsewhere is killed by every matrix entry of that minimal
differential.  It is a nonzero kernel vector.  Thus $N\leq1$,
which proves the asserted tensor-homology vanishing.
\end{proof}

\begin{theorem}[The two-column graph localization]
\label{C:two-column-graph-localization}
Let $M$ be a finite monoid with minimal rectangular band
$K=I\times\{0,1\}$, put $R=\mathbb QM$ and
$A=R/\mathbb QK$, and write
$V=\widetilde{\mathbb QI}$ for the left row augmentation module.
Let $M$ act on the right of a nonempty finite graph $\Gamma$,
mapping vertices to vertices and every edge at unit speed to an
edge, in either orientation, or to a vertex.  Set
\[
 X=\operatorname{hocofib}(\Gamma_{hM}\to BM).
\]
Either the action has a common fixed vertex, or there is an actual
two-term complex $D$ of finite-dimensional right $A$-modules such
that, naturally for every $n\geq4$,
\[
 H_n(X;\mathbb Q)
   \cong H_{n-2}(D\otimes_A^{\mathbb L}V).
\]
In particular:
\begin{enumerate}
\item If $V$ is projective over $A$, then all rational degree-two
products on $X$ vanish, without a bounded-homology hypothesis.
\item If $|I|=2$ and $A$ is local, then bounded rational homology
of $X$ implies that all its rational degree-two products vanish.
The same conclusion holds when the one-dimensional $V$ is
supported on a local direct-factor algebra of $A$.
\end{enumerate}
\end{theorem}
\begin{proof}
Choose $i_0\in I$ and put $e=x_{i_0,0}$ and $e_1=x_{i_0,1}$.
The image graphs $\Gamma_0=\Gamma e$ and $\Gamma_1=\Gamma e_1$
are retract subgraphs.  The image of any $x_{ij}$ depends only
on its column, since the same-column idempotents fix each other's
images.  Right multiplication by $e_1$ and $e$ gives mutually
inverse graph isomorphisms between $\Gamma_0$ and $\Gamma_1$.

If a vertex $u$ belongs to both image graphs, it is fixed by
both $e,e_1$ and hence by every band element:
$ux_{ij}=(ue)x_{ij}=ux_{i_0,j}=u$.
For arbitrary $m\in M$, the equality
$um=(ue)m=u(em)=u$ then makes it a common fixed vertex.
Its action-bar section makes the projection injective on
cohomology.  The cofiber exact sequence therefore makes every
positive-degree cohomology class of the cofiber a pullback
from the suspension of the action bar.  Products of such
classes vanish, as in Theorem~\ref{C:P5-graph-cofiber-cups}.

Otherwise $\Gamma_0$ and $\Gamma_1$ have disjoint vertex sets
and therefore disjoint edge sets.  Let
\[
 C=[C_1(\Gamma;\mathbb Q)\longrightarrow
       \widetilde C_0(\Gamma;\mathbb Q)]
\]
be its augmented graph complex of right $R$-modules.  Multiplication
gives a chain map
\[
 C e\otimes_{\mathbb Q}eR\longrightarrow C.
\]
This is degreewise injective.  In degree one its two column
summands have disjoint edge supports.  In degree zero,
$C_0e=\widetilde{\mathbb Q(\operatorname{Vert}\Gamma_0)}$;
the two summands have disjoint vertex supports and each has
augmentation zero.  These are precisely the image subspaces of
the two column copies.  The same argument allows reversed edge
orientations, since their only effect on chain coordinates is a
sign.

The image is exactly $C\mathbb QK$, by the literal band
factorization $\mathbb QK=Re\otimes_{\mathbb Q}eR$.
Consequently the quotient
\[
 D=C/(Ce\otimes_{\mathbb Q}eR)
\]
is literally a two-term right $A$-complex.  The subcomplex has
projective right $R$-modules in its two degrees.  After derived
tensoring with the trivial left $R$-module, it becomes the finite
complex $Ce$ in those same degrees, because $eR\otimes_R\mathbb Q
\cong\mathbb Q$.  Thus the short exact sequence of complexes
induces
\[
 H_j(C\otimes_R^{\mathbb L}\mathbb Q)
 \cong H_j(D\otimes_R^{\mathbb L}\mathbb Q),
 \qquad j\geq3.
\]

The row augmentation sequence is
\[
 0\longrightarrow V\longrightarrow Re\longrightarrow\mathbb Q
       \longrightarrow0.
\]
The ideal $\mathbb QK$ annihilates $V$, and $De=0$.
Since $Re$ is projective, tensoring this sequence gives
\[
 D\otimes_R^{\mathbb L}\mathbb Q
       \simeq (D\otimes_R^{\mathbb L}V)[1].
\]
The comparison to $A$ has no omitted derived correction.
Indeed $0\to\mathbb QK\to R\to A\to0$ is a left-projective
resolution of $A$, and a right $A$-module $U$ satisfies
$U\otimes_R\mathbb QK=Ue\otimes eR=0$.  Hence
$U\otimes_R^{\mathbb L}A\simeq U$ and
$U\otimes_R^{\mathbb L}V\simeq U\otimes_A^{\mathbb L}V$.
The same natural comparison applies to the finite complex $D$.
Finally the graph-cofiber chain model is
$(C\otimes_R^{\mathbb L}\mathbb Q)[1]$ in reduced homology.
Combining the two shifts proves the displayed formula exactly
in the asserted range $n\geq4$.

If $V$ is $A$-projective, its tensor with $D$ is concentrated
in degrees zero and one, so $H_4(X;\mathbb Q)=0$.
If $|I|=2$, the module $V$ is one-dimensional.  When $A$ is
local, its action on $V$ identifies the residue division algebra
with $\mathbb Q$, and $V$ is the residue module.  Bounded rational
homology of $X$ makes $D\otimes_A^{\mathbb L}V$ homologically
bounded.  Lemma~\ref{C:local-two-term-derived-tensor} gives
$H_2(D\otimes_A^{\mathbb L}V)=0$, so again $H_4(X;\mathbb Q)=0$.
For a local direct factor $B=Az$, with $z$ central and $zV=V$,
replace $D$ by $Dz$ and $A$ by $B$ in the same argument.
This proves all assertions, including the fixed-vertex alternative.
\end{proof}

\begin{corollary}[All sources of order at most six opposite $R_3$]
\label{C:six-source-right-zero-all-states}
Let $R_3$ be the identity-adjoined two-element right-zero band.
For every finite $(R_3,M)$-biset with $|M|\leq6$, bounded rational
homology of its collage implies that all rational degree-two
products vanish.  In particular no such collage is a projective
plane, at any number of cross states.
\end{corollary}
\begin{proof}
The left action category has the natural equivariant finite graph
model, and $BR_3$ is contractible.  Thus the collage is precisely
the graph cofiber in the theorem.
If $H^2(BM;\mathbb Q)=0$, every degree-two class is relative to
the two vertex spaces and the suspension quotient kills products.
Otherwise the minimal ideal must have at least two rows and two
columns.  At order at most six it is a $2\times2$ band with at
most two outside elements; a nontrivial local group would already
make the ideal have at least eight elements.

The contracted algebra is either $\mathbb Q$, or has basis
$1,a$ with $a^2$ equal to $0,a$, or $1$.  In the last two cases
it is semisimple, so the row module is projective.  In the first
it is the local dual-number algebra.  The preceding theorem
therefore applies in every case.  This argument uses bounded
rational homology only in the local-algebra alternative, and
does not infer finite CW type for an intermediate graph space.
\end{proof}
